% ==============================================================================
% CHAPITRE 11 : THÉORÈME DE THALÈS & AGRANDISSEMENTS / RÉDUCTIONS (VERSION ÉLÈVE)
% ==============================================================================
\chapter{Théorème de Thalès \& Agrandissements -- Réductions}

% ------------------------------------------------------------------------------
% 1. ACTIVITÉ D'INVESTIGATION (SITUATION PROFESSIONNELLE : MENUISERIE / BTP)
% ------------------------------------------------------------------------------
\begin{activite}{Sécurité et Stabilité d'un Escabeau de Chantier (Menuiserie / BTP)}
Un artisan menuisier fabrique un escabeau double en bois composé de deux montants $[AB]$ et $[AC]$ articulés au sommet $A$.
Pour garantir la stabilité et respecter les normes de sécurité de l'atelier, une traverse de renfort horizontale $[MN]$ est fixée entre les deux montants, parallèlement au sol horizontal $[BC]$.

\vspace{1mm}
\begin{center}
\begin{tcolorbox}[colback=amber!10, colframe=amber!80!black, arc=4pt, boxrule=0.8pt, left=8pt, right=8pt, top=4pt, bottom=4pt]
\sffamily\footnotesize
\textbf{Données techniques relevées sur le plan d'atelier :}
\begin{itemize}[leftmargin=1.5em, itemsep=1pt]
    \item Longueur totale du montant : $AB = 150\text{ cm}$ ;
    \item Position de la fixation de la traverse : $AM = 60\text{ cm}$ (donc $M$ appartient à $[AB]$) ;
    \item Écartement au sol des pieds : $BC = 80\text{ cm}$ ;
    \item La traverse $(MN)$ est strictement parallèle à la ligne de sol $(BC)$.
\end{itemize}
\end{tcolorbox}
\end{center}

\vspace{1mm}
\noindent
\begin{minipage}[c]{0.48\linewidth}
\centering
\begin{tikzpicture}[scale=0.95]
    % Sol horizontal
    \draw[slate600, thick] (-0.5,0) -- (4.5,0);
    \fill[pattern=north east lines, pattern color=slate300] (-0.5,-0.2) rectangle (4.5,0);
    
    % Montants escabeau
    \coordinate (A) at (2.0,3.6);
    \coordinate (B) at (0.4,0);
    \coordinate (C) at (3.6,0);
    
    % Points M et N (à 40% depuis A)
    \coordinate (M) at ($(A)!0.4!(B)$);
    \coordinate (N) at ($(A)!0.4!(C)$);
    
    % Dessin montants
    \draw[slate800, line width=2pt] (A) -- (B);
    \draw[slate800, line width=2pt] (A) -- (C);
    
    % Traverse de renfort (MN)
    \draw[colPrimary, line width=2.2pt] (M) -- (N) node[midway, above=1pt, font=\bfseries\scriptsize, text=colPrimary] {Traverse ($MN$)};
    
    % Écartement sol
    \draw[dashed, slate600, thick] (B) -- (C) node[midway, below=2pt, font=\bfseries\scriptsize] {$BC = 80\text{ cm}$};
    
    % Cotes et labels
    \node[above=2pt, font=\bfseries\small] at (A) {$A$};
    \node[below left=1pt, font=\bfseries\small] at (B) {$B$};
    \node[below right=1pt, font=\bfseries\small] at (C) {$C$};
    \node[left=2pt, font=\bfseries\small, text=colSecondary] at (M) {$M$};
    \node[right=2pt, font=\bfseries\small, text=colSecondary] at (N) {$N$};
    
    % Flèches de cote
    \draw[<->, colSecondary, thick] ($(A)+(-0.35,0)$) -- ($(M)+(-0.35,0)$) node[midway, left=1pt, font=\bfseries\tiny] {$60\text{ cm}$};
    \draw[<->, slate600, thick] ($(A)+(-0.8,0)$) -- ($(B)+(-0.8,0)$) node[midway, left=1pt, font=\bfseries\tiny] {$150\text{ cm}$};
\end{tikzpicture}
\end{minipage}\hfill
\begin{minipage}[c]{0.48\linewidth}
\begin{tcolorbox}[colback=slate100, colframe=colPrimary, arc=5pt, boxrule=0.8pt, top=6pt, bottom=6pt, left=8pt, right=8pt]
\sffamily\small
\textbf{Problématique :}
\begin{enumerate}[leftmargin=1.2em, itemsep=4pt]
    \item Les triangles $AMN$ et $ABC$ forment-ils une configuration de proportionnalité ?
    \item Calculer le rapport de réduction $k = \frac{AM}{AB}$ sous forme décimale.
    \item À l'aide de l'égalité de Thalès $\frac{AM}{AB} = \frac{MN}{BC}$, calculer la longueur exacte de la traverse $[MN]$ à couper.
\end{enumerate}
\end{tcolorbox}
\end{minipage}

\vspace{3mm}
\lignesreponse[6]
\end{activite}

\newpage

% ------------------------------------------------------------------------------
% 2. COURS : THÉORÈME DE THALÈS & PROPORTIONNALITÉ
% ------------------------------------------------------------------------------
\section{Théorème de Thalès : Calcul d'une Longueur}

\begin{cadredef}[Les 2 Configurations Clés de Thalès]
\noindent
\begin{minipage}[c]{0.52\linewidth}
Soient deux droites sécantes en $A$, coupées par deux droites $(MN)$ et $(BC)$.
\textbf{Si $(MN) \parallel (BC)$, alors :}
\[ \mathbf{\frac{AM}{AB} = \frac{AN}{AC} = \frac{MN}{BC}} \]
\end{minipage}\hfill
\begin{minipage}[c]{0.46\linewidth}
\centering
\begin{tikzpicture}[scale=0.58]
    % Config 1 : Imbriqué
    \node[above=2pt, font=\bfseries\tiny, text=colPrimary] at (1.2,2.1) {1. Imbriqué};
    \coordinate (A) at (1.2,1.9);
    \coordinate (B) at (0,0.3);
    \coordinate (C) at (2.4,0.3);
    \coordinate (M) at ($(A)!0.55!(B)$);
    \coordinate (N) at ($(A)!0.55!(C)$);
    \draw[slate800, line width=1pt] (A) -- (B) -- (C) -- cycle;
    \draw[colPrimary, line width=1.4pt] (M) -- (N);
    \node[above=1pt, font=\bfseries\tiny] at (A) {$A$};
    \node[below left=1pt, font=\bfseries\tiny] at (B) {$B$};
    \node[below right=1pt, font=\bfseries\tiny] at (C) {$C$};
    \node[left=1pt, font=\bfseries\tiny] at (M) {$M$};
    \node[right=1pt, font=\bfseries\tiny] at (N) {$N$};

    % Config 2 : Sablier
    \node[above=2pt, font=\bfseries\tiny, text=colPrimary] at (4.7,2.1) {2. Sablier};
    \coordinate (A2) at (4.7,1.1);
    \coordinate (B2) at (3.7,1.9);
    \coordinate (C2) at (5.7,1.9);
    \coordinate (M2) at (3.9,0.3);
    \coordinate (N2) at (5.5,0.3);
    \draw[slate800, line width=1pt] (B2) -- (N2);
    \draw[slate800, line width=1pt] (C2) -- (M2);
    \draw[colPrimary, line width=1.4pt] (B2) -- (C2);
    \draw[colPrimary, line width=1.4pt] (M2) -- (N2);
    \node[right=2pt, font=\bfseries\tiny] at (A2) {$A$};
    \node[above left=1pt, font=\bfseries\tiny] at (B2) {$B$};
    \node[above right=1pt, font=\bfseries\tiny] at (C2) {$C$};
    \node[below left=1pt, font=\bfseries\tiny] at (M2) {$M$};
    \node[below right=1pt, font=\bfseries\tiny] at (N2) {$N$};
\end{tikzpicture}
\end{minipage}
\end{cadredef}

\begin{cadremethode}[Exemple Résolu : Calcul d'une Longueur par Produit en Croix]
\noindent
\begin{minipage}[c]{0.62\linewidth}
\textbf{Énoncé} : Dans le triangle $ABC$, $(DE) \parallel (BC)$ avec $AB = 10\text{ cm}$, $AD = 4\text{ cm}$ et $BC = 15\text{ cm}$. Calculer la longueur $DE$.
\begin{enumerate}[leftmargin=1.2em, itemsep=1pt]
    \item Les droites $(DB)$ et $(EC)$ sont sécantes en $A$, et $(DE) \parallel (BC)$.
    \item D'après le théorème de Thalès : $\frac{AD}{AB} = \frac{DE}{BC} \implies \frac{4}{10} = \frac{DE}{15}$.
    \item Produit en croix : $DE = \frac{4 \times 15}{10} = \frac{60}{10} = \mathbf{6\text{ cm}}$.
\end{enumerate}
\end{minipage}\hfill
\begin{minipage}[c]{0.36\linewidth}
\centering
\begin{tikzpicture}[scale=0.55]
    \coordinate (A) at (1.5,2.4);
    \coordinate (B) at (0,0);
    \coordinate (C) at (3.0,0);
    \coordinate (D) at ($(A)!0.4!(B)$);
    \coordinate (E) at ($(A)!0.4!(C)$);
    \draw[slate800, line width=1.2pt] (A) -- (B) -- (C) -- cycle;
    \draw[colPrimary, line width=1.6pt] (D) -- (E) node[midway, above=1pt, font=\bfseries\tiny] {$DE = ?$};
    \node[above, font=\bfseries\scriptsize] at (A) {$A$};
    \node[below left, font=\bfseries\scriptsize] at (B) {$B$};
    \node[below right, font=\bfseries\scriptsize] at (C) {$C$};
    \node[left, font=\bfseries\scriptsize] at (D) {$D$};
    \node[right, font=\bfseries\scriptsize] at (E) {$E$};
    \node[below=2pt, font=\bfseries\tiny] at (1.5,0) {$15\text{ cm}$};
    \node[left=2pt, font=\bfseries\tiny] at ($(A)!0.2!(D)$) {$4\text{ cm}$};
    \node[left=2pt, font=\bfseries\tiny] at ($(A)!0.7!(B)$) {$10\text{ cm}$};
\end{tikzpicture}
\end{minipage}
\end{cadremethode}

% ------------------------------------------------------------------------------
% 3. COURS : AGRANDISSEMENTS & RÉDUCTIONS (LONGUEURS, AIRES, VOLUMES)
% ------------------------------------------------------------------------------
\section{Agrandissements \& Réductions : Effet sur $k$, $k^2$, $k^3$}

\begin{cadredef}[Représentation Graphique de l'Effet du Coefficient $k$]
\noindent
\begin{minipage}[c]{0.45\linewidth}
Lorsqu'on multiplie les longueurs par $\mathbf{k}$ ($k > 0$) :
\begin{itemize}[leftmargin=1.2em, itemsep=1pt]
    \item \textbf{Longueurs} $\implies \mathbf{\times k}$
    \item \textbf{Aires} $\implies \mathbf{\times k^2}$
    \item \textbf{Volumes} $\implies \mathbf{\times k^3}$
\end{itemize}
\end{minipage}\hfill
\begin{minipage}[c]{0.53\linewidth}
\centering
\begin{tikzpicture}[scale=0.52]
    % Cube 1 (k=1)
    \draw[slate800, fill=colPrimary!20, thick] (0,0) rectangle (1.2,1.2);
    \draw[slate800, fill=colPrimary!30, thick] (1.2,0) -- (1.6,0.4) -- (1.6,1.6) -- (1.2,1.2) -- cycle;
    \draw[slate800, fill=colPrimary!10, thick] (0,1.2) -- (0.4,1.6) -- (1.6,1.6) -- (1.2,1.2) -- cycle;
    \node[below=2pt, font=\bfseries\tiny] at (0.6,0) {$c = 1$};
    \node[below=7pt, font=\bfseries\tiny, align=center] at (0.8,-0.05) {$\text{Aire} = 1$\\[-1pt]$\text{Vol} = 1$};

    % Flèche agrandissement k=2
    \draw[->, colSecondary, line width=1.5pt] (1.9,0.8) -- (2.9,0.8) node[midway, above=1pt, font=\bfseries\tiny] {$\times k = 2$};

    % Cube 2 (k=2) avec 8 cubes visibles
    \draw[slate800, fill=colSecondary!20, thick] (3.3,0) rectangle (5.7,2.4);
    \draw[slate800, fill=colSecondary!30, thick] (5.7,0) -- (6.5,0.8) -- (6.5,3.2) -- (5.7,2.4) -- cycle;
    \draw[slate800, fill=colSecondary!10, thick] (3.3,2.4) -- (4.1,3.2) -- (6.5,3.2) -- (5.7,2.4) -- cycle;
    \draw[slate600, dashed] (4.5,0) -- (4.5,2.4);
    \draw[slate600, dashed] (3.3,1.2) -- (5.7,1.2);
    \draw[slate600, dashed] (5.7,1.2) -- (6.5,2.0);
    \draw[slate600, dashed] (4.5,2.4) -- (5.3,3.2);
    \node[below=2pt, font=\bfseries\tiny] at (4.5,0) {$c' = 2\ (\mathbf{\times 2})$};
    \node[below=7pt, font=\bfseries\tiny, align=center] at (4.9,-0.05) {$\text{Aire}' = 4\ (\mathbf{\times 2^2})$\\[-1pt]$\text{Vol}' = 8\ (\mathbf{\times 2^3})$};
\end{tikzpicture}
\end{minipage}
\end{cadredef}

\begin{cadreexemple}[Exemple Chiffré : Agrandissement d'une Boîte Industrielle]
Un conteneur a une longueur de $L = 5\text{ m}$, une surface de base de $S = 12\text{ m}^2$ et un volume de $V = 30\text{ m}^3$.
On applique un coefficient d'agrandissement $\mathbf{k = 3}$ :
\begin{itemize}[leftmargin=1.5em, itemsep=1pt]
    \item \textbf{Nouvelle longueur} : $L' = L \times k = 5 \times 3 = \mathbf{15\text{ m}}$ ;
    \item \textbf{Nouvelle surface} : $S' = S \times k^2 = 12 \times 3^2 = 12 \times 9 = \mathbf{108\text{ m}^2}$ ;
    \item \textbf{Nouveau volume} : $V' = V \times k^3 = 30 \times 3^3 = 30 \times 27 = \mathbf{810\text{ m}^3}$.
\end{itemize}
\end{cadreexemple}

\newpage

% ------------------------------------------------------------------------------
% 4. QUESTIONS FLASH / AUTOMATISMES & TD 1 & TD 2 (PAGE 3)
% ------------------------------------------------------------------------------
\section{Questions Flash / Automatismes}

\begin{automatisme}[8 Questions Flash d'Entraînement]
\begin{enumerate}[label=\textbf{Q\arabic*.}, itemsep=1pt]
    \item Configuration de Thalès avec $(MN) \parallel (BC)$ : $\frac{AM}{AB} = \frac{\pointille[2cm]}{AC} = \frac{MN}{\pointille[2cm]}$
    \item Si $\frac{x}{6} = \frac{5}{10}$, calculer la valeur exacte de $x$ : $x = \pointille[2cm]$
    \item Échelle $1:50$ : s'agit-il d'un agrandissement ou d'une réduction ? \pointille[3cm]
    \item Un carré de côté $4\text{ cm}$ est agrandi avec $k = 3$. Nouveau côté : \pointille[2cm]
    \item Si les longueurs sont multipliées par $k = 2$, l'aire est multipliée par : \pointille[2cm]
    \item Volume initial $100\text{ cm}^3$, agrandi avec $k=2$. Nouveau volume : \pointille[2cm]
    \item Donner sous forme décimale le rapport $k = \frac{3}{12} = \pointille[2cm] $
    \item Droites parallèles si les quotients de Thalès sont \pointille[3cm]
\end{enumerate}
\end{automatisme}

\section{Travaux Dirigés (TD) : Problèmes d'Entraînement}

\begin{exercice}[2 pts][\capRealiser]{Renfort d'une Ferme de Toit (Charpente / BTP)}
\noindent
\begin{minipage}[c]{0.58\linewidth}
Une ferme de charpente $ABC$ est renforcée par un entrait horizontal $[MN]$ parallèle à $[BC]$.
On donne : $AB = 6\text{ m}$, $AM = 4\text{ m}$, et $BC = 7{,}50\text{ m}$.
\begin{enumerate}[leftmargin=1.2em, itemsep=1pt]
    \item Calculer la longueur de la poutre $[MN]$.
    \item Sachant que $AC = 5{,}40\text{ m}$, calculer $AN$.
\end{enumerate}
\end{minipage}\hfill
\begin{minipage}[c]{0.38\linewidth}
\centering
\begin{tikzpicture}[scale=0.62]
    \coordinate (A) at (2.0,2.5);
    \coordinate (B) at (0,0);
    \coordinate (C) at (4.0,0);
    \coordinate (M) at ($(A)!0.666!(B)$);
    \coordinate (N) at ($(A)!0.666!(C)$);
    \draw[slate800, line width=1.4pt] (A) -- (B) -- (C) -- cycle;
    \draw[colPrimary, line width=1.8pt] (M) -- (N) node[midway, above=1pt, font=\bfseries\tiny] {$MN$};
    \node[above=1pt, font=\bfseries\scriptsize] at (A) {$A$};
    \node[below left=1pt, font=\bfseries\scriptsize] at (B) {$B$};
    \node[below right=1pt, font=\bfseries\scriptsize] at (C) {$C$};
    \node[left=1pt, font=\bfseries\scriptsize] at (M) {$M$};
    \node[right=1pt, font=\bfseries\scriptsize] at (N) {$N$};
\end{tikzpicture}
\end{minipage}

\vspace{1mm}
\lignesreponse[2]
\end{exercice}

\begin{exercice}[3 pts][\capSApproprier]{Sangle de Sécurité d'un Escabeau Professionnel}
\noindent
\begin{minipage}[c]{0.58\linewidth}
Un escabeau de peintre possède une sangle $[DE]$ parallèle au sol $[BC]$.
On a : $AD = 45\text{ cm}$, $AB = 120\text{ cm}$, et $DE = 30\text{ cm}$.
\begin{enumerate}[leftmargin=1.2em, itemsep=1pt]
    \item Calculer l'écartement au sol $BC$.
    \item Si le montant mesure $AC = 120\text{ cm}$, calculer $AE$.
\end{enumerate}
\end{minipage}\hfill
\begin{minipage}[c]{0.38\linewidth}
\centering
\begin{tikzpicture}[scale=0.62]
    \coordinate (A) at (1.8,2.6);
    \coordinate (B) at (0.3,0);
    \coordinate (C) at (3.3,0);
    \coordinate (D) at ($(A)!0.375!(B)$);
    \coordinate (E) at ($(A)!0.375!(C)$);
    \draw[slate600, thick] (0,0) -- (3.6,0);
    \fill[pattern=north east lines, pattern color=slate300] (0,-0.2) rectangle (3.6,0);
    \draw[slate800, line width=1.4pt] (A) -- (B);
    \draw[slate800, line width=1.4pt] (A) -- (C);
    \draw[colSecondary, line width=1.8pt] (D) -- (E) node[midway, above=1pt, font=\bfseries\tiny] {$30\text{ cm}$};
    \node[above=1pt, font=\bfseries\scriptsize] at (A) {$A$};
    \node[below left=1pt, font=\bfseries\scriptsize] at (B) {$B$};
    \node[below right=1pt, font=\bfseries\scriptsize] at (C) {$C$};
    \node[left=1pt, font=\bfseries\scriptsize] at (D) {$D$};
    \node[right=1pt, font=\bfseries\scriptsize] at (E) {$E$};
\end{tikzpicture}
\end{minipage}

\vspace{1mm}
\lignesreponse[2]
\end{exercice}

\newpage

% ------------------------------------------------------------------------------
% 5. TRAVAUX DIRIGÉS (TD) - EXERCICES 3, 4 & 5 (PAGE 4)
% ------------------------------------------------------------------------------
\begin{exercice}[3 pts][\capAnalyser]{Mesure de la Largeur d'une Rivière (Configuration Sablier)}
\noindent
\begin{minipage}[c]{0.58\linewidth}
Un géomètre mesure la distance $AB$ séparant les deux rives d'un fleuve.
Les droites $(AD)$ et $(BC)$ se coupent en $O$, avec $(AB) \parallel (CD)$.
Mesures : $OC = 18\text{ m}$, $OB = 45\text{ m}$, et $CD = 24\text{ m}$.
\begin{enumerate}[leftmargin=1.2em, itemsep=1pt]
    \item Écrire l'égalité des trois rapports de Thalès en $O$.
    \item Calculer la largeur de la rivière $AB$.
\end{enumerate}
\end{minipage}\hfill
\begin{minipage}[c]{0.38\linewidth}
\centering
\begin{tikzpicture}[scale=0.62]
    \fill[sky!15] (-0.3,0.7) rectangle (4.2,1.8);
    \draw[sky!80!blue, dashed] (-0.3,0.7) -- (4.2,0.7);
    \draw[sky!80!blue, dashed] (-0.3,1.8) -- (4.2,1.8);
    \coordinate (A) at (0.4,2.3);
    \coordinate (B) at (3.2,2.3);
    \coordinate (O) at (1.8,1.25);
    \coordinate (C) at (1.0,0.2);
    \coordinate (D) at (2.6,0.2);
    \draw[colPrimary, line width=1.4pt] (A) -- (D);
    \draw[colPrimary, line width=1.4pt] (B) -- (C);
    \draw[slate800, line width=1.6pt] (A) -- (B) node[midway, above=1pt, font=\bfseries\tiny] {$AB$};
    \draw[slate800, line width=1.6pt] (C) -- (D) node[midway, below=1pt, font=\bfseries\tiny] {$24\text{ m}$};
    \node[above left=1pt, font=\bfseries\scriptsize] at (A) {$A$};
    \node[above right=1pt, font=\bfseries\scriptsize] at (B) {$B$};
    \node[right=1pt, font=\bfseries\scriptsize] at (O) {$O$};
    \node[below left=1pt, font=\bfseries\scriptsize] at (C) {$C$};
    \node[below right=1pt, font=\bfseries\scriptsize] at (D) {$D$};
\end{tikzpicture}
\end{minipage}

\vspace{1mm}
\lignesreponse[2]
\end{exercice}

\begin{exercice}[3 pts][\capValider]{Contrôle d'Horizontalité d'une Étagère (Réciproque de Thalès)}
\noindent
\begin{minipage}[c]{0.58\linewidth}
Un poseur vérifie si une tablette d'étagère $[MN]$ est bien parallèle au montant bas $[BC]$.
Il mesure : $AM = 24\text{ cm}$, $AB = 60\text{ cm}$, $AN = 20\text{ cm}$, et $AC = 50\text{ cm}$.
\begin{enumerate}[leftmargin=1.2em, itemsep=1pt]
    \item Calculer et comparer les quotients $\frac{AM}{AB}$ et $\frac{AN}{AC}$.
    \item Conclure sur le parallélisme de la tablette $(MN)$.
\end{enumerate}
\end{minipage}\hfill
\begin{minipage}[c]{0.38\linewidth}
\centering
\begin{tikzpicture}[scale=0.62]
    \coordinate (A) at (1.8,2.4);
    \coordinate (B) at (0.2,0);
    \coordinate (C) at (3.4,0);
    \coordinate (M) at ($(A)!0.4!(B)$);
    \coordinate (N) at ($(A)!0.4!(C)$);
    \draw[slate800, line width=1.4pt] (A) -- (B) -- (C) -- cycle;
    \draw[colSecondary, line width=1.8pt] (M) -- (N) node[midway, above=1pt, font=\bfseries\tiny] {Tablette ($MN$)};
    \node[above=1pt, font=\bfseries\scriptsize] at (A) {$A$};
    \node[below left=1pt, font=\bfseries\scriptsize] at (B) {$B$};
    \node[below right=1pt, font=\bfseries\scriptsize] at (C) {$C$};
    \node[left=1pt, font=\bfseries\scriptsize] at (M) {$M$};
    \node[right=1pt, font=\bfseries\scriptsize] at (N) {$N$};
\end{tikzpicture}
\end{minipage}

\vspace{1mm}
\lignesreponse[2]
\end{exercice}

\begin{exercice}[3 pts][\capRealiser]{Maquette d'Architecture et Surface d'un Atelier}
\noindent
\begin{minipage}[c]{0.58\linewidth}
Un architecte réalise la maquette à l'échelle $\mathbf{1:50}$ ($k = 0{,}02$) d'un atelier rectangulaire de dimensions réelles $L = 25\text{ m}$ et $l = 12\text{ m}$.
\begin{enumerate}[leftmargin=1.2em, itemsep=1pt]
    \item Calculer les dimensions de la maquette $L'$ et $l'$ en centimètres.
    \item L'aire réelle vaut $300\text{ m}^2$. Calculer l'aire de la maquette en $\text{cm}^2$.
\end{enumerate}
\end{minipage}\hfill
\begin{minipage}[c]{0.38\linewidth}
\centering
\begin{tikzpicture}[scale=0.62]
    \draw[colPrimary, fill=colPrimary!10, line width=1.1pt] (0,0.6) rectangle (2.0,1.9);
    \node[font=\bfseries\tiny] at (1.0,1.25) {Atelier Réel};
    \node[below, font=\bfseries\tiny] at (1.0,0.6) {$25\text{ m}$};
    \draw[colSecondary, fill=colSecondary!15, line width=1.1pt] (2.5,0.6) rectangle (3.4,1.1);
    \node[font=\bfseries\tiny] at (2.95,0.85) {Maquette};
\end{tikzpicture}
\end{minipage}

\vspace{1mm}
\lignesreponse[2]
\end{exercice}

\newpage

% ------------------------------------------------------------------------------
% 5. TRAVAUX DIRIGÉS (TD) & TP TICE (PAGE 5)
% ------------------------------------------------------------------------------
\begin{exercice}[3 pts][\capAnalyser]{Hauteur d'un Bâtiment par l'Ombre Portée (Méthode de Thalès)}
\noindent
\begin{minipage}[c]{0.58\linewidth}
Les rayons du soleil sont considérés comme parallèles.
Un piquet vertical de $1{,}50\text{ m}$ projette une ombre au sol de $2{,}00\text{ m}$.
Au même moment, l'ombre d'un bâtiment mesure $36{,}00\text{ m}$.
\begin{enumerate}[leftmargin=1.2em, itemsep=1pt]
    \item Justifier pourquoi les triangles d'ombre sont semblables.
    \item Calculer la hauteur réelle $H$ du bâtiment.
\end{enumerate}
\end{minipage}\hfill
\begin{minipage}[c]{0.38\linewidth}
\centering
\begin{tikzpicture}[scale=0.62]
    \draw[slate600, thick] (-0.3,0) -- (4.0,0);
    \fill[pattern=north east lines, pattern color=slate300] (-0.3,-0.2) rectangle (4.0,0);
    \draw[slate800, fill=slate200, line width=1.1pt] (0,0) rectangle (1.1,1.9);
    \node[font=\bfseries\tiny] at (0.55,0.95) {Bâtiment};
    \draw[amber!90!black, line width=1.2pt, dashed] (1.1,1.9) -- (3.3,0);
    \node[below=1pt, font=\bfseries\tiny] at (2.2,0) {Ombre $= 36\text{ m}$};
    \draw[colPrimary, line width=1.6pt] (3.5,0) -- (3.5,0.45);
\end{tikzpicture}
\end{minipage}

\vspace{1mm}
\lignesreponse[2]
\end{exercice}

\begin{exercice}[4 pts][\capRealiser]{Impression 3D \& Réduction de Volume d'un Flacon}
\noindent
\begin{minipage}[c]{0.58\linewidth}
Une entreprise de cosmétique fabrique un flacon de contenance $V = 250\text{ mL}$ et de hauteur $h = 15\text{ cm}$.
Elle crée un échantillon miniature avec un coefficient $\mathbf{k = 0{,}6}$.
\begin{enumerate}[leftmargin=1.2em, itemsep=1pt]
    \item Calculer la hauteur $h'$ de l'échantillon miniature.
    \item Calculer la contenance $V'$ de l'échantillon (en $\text{mL}$, rappel : $\times k^3$).
\end{enumerate}
\end{minipage}\hfill
\begin{minipage}[c]{0.38\linewidth}
\centering
\begin{tikzpicture}[scale=0.62]
    \draw[colPrimary, fill=colPrimary!10, line width=1.1pt] (0.4,0) -- (1.4,0) -- (1.4,1.6) -- (1.1,1.9) -- (1.1,2.2) -- (0.7,2.2) -- (0.7,1.9) -- (0.4,1.6) -- cycle;
    \node[font=\bfseries\tiny] at (0.9,0.8) {$250\text{ mL}$};
    \draw[colSecondary, fill=colSecondary!15, line width=1.1pt] (2.4,0) -- (3.0,0) -- (3.0,0.96) -- (2.82,1.14) -- (2.82,1.32) -- (2.58,1.32) -- (2.58,1.14) -- (2.4,0.96) -- cycle;
    \node[font=\bfseries\tiny] at (2.7,0.48) {$V'$};
\end{tikzpicture}
\end{minipage}

\vspace{1mm}
\lignesreponse[2]
\end{exercice}

% ------------------------------------------------------------------------------
% 6. TP TICE / CALCULATRICE NUMWORKS
% ------------------------------------------------------------------------------
\section{TP TICE : Proportions \& Équations sur NumWorks}

\begin{tpinfo}[Calculatrice NumWorks]{Résolution d'Égalités de Rapports \& Fractions}
\begin{enumerate}[leftmargin=1.5em, itemsep=2pt]
    \item \textbf{Calcul direct d'une quatrième proportionnelle (Application Calculs)} :
    \begin{itemize}
        \item Pour résoudre $\frac{x}{15} = \frac{4}{10}$, taper directement \texttt{(4 * 15) / 10} puis valider avec \textbf{EXE}.
        \item Utiliser la touche fraction $\mathbf{\frac{\Box}{\Box}}$ pour saisir directement les écritures rationnelles.
    \end{itemize}
    \item \textbf{Résolution automatique d'équations (Application Équations)} :
    \begin{itemize}
        \item Ouvrir l'application \textbf{Équations} $\rightarrow$ \textbf{Ajouter une équation}.
        \item Saisir : \texttt{x/15 = 4/10} puis sélectionner \textbf{Résoudre l'équation}.
    \end{itemize}
\end{enumerate}
\end{tpinfo}

\newpage

% ------------------------------------------------------------------------------
% 7. VERS LE BREVET (DNB PRO) (PAGE 6)
% ------------------------------------------------------------------------------
\section{Vers le Brevet (DNB Pro)}

\begin{sujetdnb}[DNB Pro Métropole][10 pts]{Installation d'un Échafaudage et Monte-Charge de Façade}
Pour ravaler la façade d'un immeuble, une entreprise de maçonnerie monte un échafaudage vertical adossé au mur $[AB]$.
Pour monter les sacs de mortier au 2\textsuperscript{ème} étage, un câble de treuil $[AC]$ de longueur \textbf{$13\text{ m}$} est fixé au sommet $A$ ($12\text{ m}$).
Une jambe de force de maintien horizontal $[MN]$ est fixée à mi-hauteur ($AM = 6\text{ m}$), parallèlement au sol horizontal $[BC]$.

\begin{center}
\begin{tikzpicture}[scale=0.82]
    % Sol
    \draw[slate600, thick] (-0.5,0) -- (6.0,0);
    \fill[pattern=north east lines, pattern color=slate300] (-0.5,-0.25) rectangle (6.0,0);

    % Façade immeuble (mur vertical AB)
    \draw[slate800, line width=2.5pt] (0.5,0) -- (0.5,3.6) node[above, font=\bfseries\small] {Sommet $A$ ($12\text{ m}$)};
    \draw[pattern=crosshatch dots, pattern color=slate300] (0,0) rectangle (0.5,3.6);
    
    % Point d'ancrage sol C
    \coordinate (A) at (0.5,3.6);
    \coordinate (B) at (0.5,0);
    \coordinate (C) at (5.3,0);
    
    % Câble de treuil (AC)
    \draw[colPrimary, line width=2pt] (A) -- (C) node[midway, above right=1pt, font=\bfseries\small] {Câble $AC = 13\text{ m}$};
    
    % Renfort horizontal (MN) à 50% de la hauteur
    \coordinate (M) at ($(A)!0.5!(B)$);
    \coordinate (N) at ($(A)!0.5!(C)$);
    \draw[colSecondary, line width=2.2pt] (M) -- (N) node[midway, above=1pt, font=\bfseries\footnotesize, text=colSecondary] {Renfort ($MN$)};
    
    % Angle droit au sol
    \draw[draw=slate800, fill=slate200, thick] (0.5,0.3) -- (0.8,0.3) -- (0.8,0) -- cycle;
    \draw[draw=slate800, fill=slate200, thick] (0.5,2.1) -- (0.8,2.1) -- (0.8,1.8) -- cycle;
    
    % Labels sommets
    \node[below left=1pt, font=\bfseries\small] at (B) {$B$};
    \node[below right=1pt, font=\bfseries\small] at (C) {$C$};
    \node[left=2pt, font=\bfseries\small] at (M) {$M$};
    \node[right=2pt, font=\bfseries\small] at (N) {$N$};
    
    % Cotes
    \node[left=2pt, font=\bfseries\footnotesize] at (0.5,2.7) {$AM = 6\text{ m}$};
    \node[left=2pt, font=\bfseries\footnotesize] at (0.5,0.9) {$MB = 6\text{ m}$};
\end{tikzpicture}
\end{center}

\begin{enumerate}[leftmargin=1.5em, itemsep=3pt]
    \item \capSApproprier\ Identifier la nature du triangle $ABC$ formé par la façade, le sol et le câble.
    \item \capRealiser\ À l'aide du théorème de Pythagore dans le triangle $ABC$ rectangle en $B$, calculer la distance au sol $BC$.
    \item \capRealiser\ Les droites $(MN)$ et $(BC)$ étant parallèles, appliquer le théorème de Thalès pour calculer la longueur du renfort $[MN]$.
    \item \capValider\ Le fabricant exige que l'aire triangulaire $AMN$ soit d'au moins \textbf{$7{,}5\text{ m}^2$} pour la stabilité. L'installation est-elle conforme ?
\end{enumerate}

\vspace{3mm}
\lignesreponse[8]
\end{sujetdnb}
